\documentclass[../../editions/book2-optimization/main.tex]{subfiles}
\begin{document}

\chapter*{Part Synthesis: Optimization Machinery}
\addcontentsline{toc}{chapter}{Part Synthesis}

\begin{keyconcept}
\textbf{What You Now Understand:}
\begin{itemize}
    \item How stochastic gradient descent trades per-step accuracy for computational efficiency
    \item How momentum and adaptive methods accelerate convergence
    \item How regularization controls model complexity and prevents overfitting
    \item How different loss functions encode different learning objectives
    \item When and why optimization converges to good solutions
\end{itemize}
\end{keyconcept}

\begin{keyconcept}
\textbf{Tools You Now Have:}
\begin{itemize}
    \item Optimizer selection: SGD, Adam, AdaGrad, RMSprop and when to use each
    \item Learning rate schedules and warmup strategies
    \item Diagnosing training: loss curves, gradient norms, validation gaps
    \item L1/L2 regularization and early stopping
\end{itemize}
\end{keyconcept}

%------------------------------------------------------------------------------
\section*{Unified Decision Tree for Optimizer Selection}

Choosing the right optimizer can feel overwhelming given the variety of options. The following decision tree provides a practical starting point based on your problem characteristics.

\begin{keyconcept}
\textbf{Optimizer Selection Flowchart:}
\begin{enumerate}
    \item \textbf{Is your problem convex?}
    \begin{itemize}
        \item Yes $\rightarrow$ Use gradient descent or L-BFGS. Convexity guarantees convergence to the global optimum with proper learning rate.
    \end{itemize}
    \item \textbf{Is it a language model or Transformer?}
    \begin{itemize}
        \item Yes $\rightarrow$ Use \textbf{AdamW with warmup}. This combination is essential for stable training of attention-based architectures.
    \end{itemize}
    \item \textbf{Is it computer vision (CNNs)?}
    \begin{itemize}
        \item Yes $\rightarrow$ Use \textbf{SGD with momentum} for best generalization, or AdamW for faster prototyping.
    \end{itemize}
    \item \textbf{Small dataset (fewer than 10,000 samples)?}
    \begin{itemize}
        \item Yes $\rightarrow$ Use \textbf{SGD with momentum} and monitor for overfitting. Adaptive methods can overfit more easily on small data.
    \end{itemize}
    \item \textbf{Is it a GAN?} $\rightarrow$ Use Adam with $\beta_1=0.5$ (lower momentum reduces oscillation between generator and discriminator)
    \item \textbf{Is it reinforcement learning?} $\rightarrow$ Use Adam (non-stationary objectives benefit from adaptivity)
    \item \textbf{Is it fine-tuning a pretrained model?} $\rightarrow$ Use AdamW with learning rate 10$\times$ lower than training from scratch, plus linear warmup
    \item \textbf{Default choice when uncertain:}
    \begin{itemize}
        \item Use \textbf{Adam} with $\eta = 0.001$, $\beta_1 = 0.9$, $\beta_2 = 0.999$. It is robust across many problem types with minimal tuning.
    \end{itemize}
\end{enumerate}
\end{keyconcept}

\begin{keyconcept}{Regularization by Architecture Type}
\begin{itemize}
\item \textbf{CNNs}: BatchNorm + weight decay ($10^{-4}$) + data augmentation + dropout (0.5 in FC layers only)
\item \textbf{Transformers}: LayerNorm (or RMSNorm) + dropout (0.1) + label smoothing (0.1) + AdamW weight decay ($10^{-2}$)
\item \textbf{RNNs/LSTMs}: LayerNorm + dropout (0.2--0.5) + gradient clipping (1.0--5.0) + weight decay ($10^{-5}$)
\end{itemize}
\end{keyconcept}

%------------------------------------------------------------------------------
\section*{Common Training Failures and Fixes}

When training goes wrong, systematic diagnosis is essential. The following table summarizes common failure modes and their remedies.

\begin{table}[htbp]
\centering
\caption{Diagnosing and fixing common training problems.}
\label{tab:training-failures}
\begin{tabular}{@{}p{3.5cm}p{5.5cm}p{5.5cm}@{}}
\toprule
\textbf{Symptom} & \textbf{Likely Causes} & \textbf{Solutions} \\
\midrule
Loss exploding (goes to infinity or NaN) & Learning rate too high; unstable gradients; numerical overflow & Lower learning rate by 10$\times$; add gradient clipping (max norm 1.0--5.0); check for division by zero \\
\addlinespace
Loss stuck / not decreasing & Learning rate too low; vanishing gradients; data pipeline bug & Increase learning rate; check data shuffling and labels; verify gradients are non-zero \\
\addlinespace
Training loss good but validation loss bad (overfitting) & Model too complex; insufficient regularization; training too long & Add dropout (0.1--0.5); increase weight decay; use early stopping; add data augmentation \\
\addlinespace
Slow convergence & Suboptimal optimizer; missing momentum; poor initialization & Switch from SGD to Adam; add momentum ($\gamma = 0.9$); use proper weight initialization \\
\addlinespace
NaN in loss or gradients & Numerical instability; log of zero or negative; exploding activations & Add epsilon guards ($10^{-8}$) to divisions; use numerically stable loss implementations; lower learning rate \\
\addlinespace
Validation loss oscillating & Learning rate too high for fine-tuning phase; batch size too small & Reduce learning rate; use learning rate scheduling; increase batch size if memory permits \\
\bottomrule
\end{tabular}
\end{table}

\begin{important}
\textbf{Debugging Priority:} When training fails, check in this order: (1) data pipeline and labels, (2) learning rate, (3) gradient magnitudes, (4) loss function implementation. Most training bugs are data bugs or learning rate issues, not optimizer problems.
\end{important}

%------------------------------------------------------------------------------
\section*{Key Takeaways}

\begin{keyconcept}
\textbf{Core Principles to Remember:}
\begin{itemize}
    \item \textbf{Convexity guarantees global optimum.} For convex problems, any local minimum is the global minimum. Non-convex optimization requires more care: good initialization, momentum to escape saddle points, and acceptance that we may find local rather than global optima.

    \item \textbf{SGD noise helps generalization.} The stochasticity in mini-batch gradients is not just a computational convenience---it acts as implicit regularization, helping models find flatter minima that generalize better. Do not over-optimize by using very large batch sizes without adjusting other hyperparameters.

    \item \textbf{Adam is robust; SGD can generalize better.} Adam converges faster and requires less tuning, making it excellent for prototyping. SGD with momentum, when properly tuned, often achieves slightly better final test performance, especially for computer vision tasks.

    \item \textbf{Regularization is about bias-variance tradeoff.} Adding regularization increases bias (the model is more constrained) but decreases variance (the model is less sensitive to training data noise). The optimal regularization strength balances these competing effects.

    \item \textbf{Learning rate is the most important hyperparameter.} If you only have time to tune one thing, tune the learning rate. It affects convergence speed, stability, and final performance more than any other single hyperparameter.

    \item \textbf{Warmup stabilizes training.} For Transformers and large batch training, gradually increasing the learning rate during early training prevents instability from unreliable initial gradient estimates.
\end{itemize}
\end{keyconcept}

%------------------------------------------------------------------------------
\section*{What Comes Next}

The optimization machinery you have learned is not just theory---it is the engine that powers every neural network training run. In Book 3, \textit{The Generative Revolution}, you will see how these optimization techniques apply to training neural networks, from simple perceptrons to modern transformers and diffusion models. You will learn how:

\begin{itemize}
    \item \textbf{Backpropagation} efficiently computes gradients through computational graphs, enabling the optimization methods you've learned to train deep networks with millions of parameters.
    \item \textbf{Automatic differentiation} frameworks (PyTorch, TensorFlow) handle gradient computation automatically, letting you focus on architecture design rather than manual calculus.
    \item \textbf{Architecture choices} interact with optimization: skip connections enable training of very deep networks by improving gradient flow; batch normalization smooths the loss landscape; attention mechanisms create different optimization challenges than convolutions.
\end{itemize}

\begin{keyconcept}
\textbf{Your Next Step:}
Book 3, \textit{The Generative Revolution}, applies these optimization techniques to train neural networks, where the architecture itself becomes a design choice.
\end{keyconcept}

\end{document}
